\documentclass[12pt]{article} % Tamanho de letra 12, papel A4, Estilo de artigo

\usepackage[utf8]{inputenc} % Pacote para permitir o uso de acentos
% Ajustes de hifenização
\tolerance=1
\emergencystretch=\maxdimen
\hyphenpenalty=10000
\hbadness=10000
\usepackage{lipsum} % Pacote para gerar texto aleatório
\usepackage{ragged2e} % Pacote para mais opções de alinhamento
\usepackage{graphicx} % Pacote para inserir figuras jpeg, pdf e png
\usepackage{ae}
\usepackage{epsfig}
\usepackage{mathtools}
\usepackage[export]{adjustbox}
\usepackage{lineno}
\usepackage{amsthm}
\usepackage{commath}
\usepackage{color, colortbl}
\usepackage{bbm}
\usepackage{setspace} % Pacote para ajustar espaçamento entre linhas

% Instructions to prevent LaTeX from hyphenating words
\tolerance=1
\emergencystretch=\maxdimen
\hyphenpenalty=10000
\hbadness=10000

\graphicspath{{figures/}} % Indica o diretório no qual as imagens estão sendo salvas, evitando que precisamos digitar parte do caminho da pasta que contém as imagens
\usepackage[rightcaption]{sidecap} % Pacote para mais opções de figura, colocar, título da legenda da figura ao lado direito
\usepackage{url} % Pacote para citar referências de sites
\usepackage{makeidx} % Pacote para criar índice
\makeindex % Habilitar a criação do índice
\usepackage{longtable} % Pacote para usar tabelas longas

% ----------
% Pacotes para dar quebra de linha dentro da tabela:
%\usepackage{fourier} 
%\usepackage{array}
\usepackage{makecell}

\renewcommand\theadalign{bc}
\renewcommand\theadfont{\bfseries}
\renewcommand\theadgape{\Gape[4pt]}
\renewcommand\cellgape{\Gape[4pt]}
% ----------

\usepackage{wrapfig} % Pacote para alinhar figuras à esquerda ou à direita do texto
\usepackage{ae}
\usepackage{epsfig}
\usepackage{mathtools}
\usepackage[export]{adjustbox}
\usepackage{lineno}
\usepackage{amsfonts}
\usepackage{amsthm}
\usepackage{commath}
\usepackage{color, colortbl}
\usepackage{bbm}

% ----------
% Pacote para inserir cabeçalho e rodapé:
% Atenção! Se remover o cabeçalho e rodapé, o LaTeX coloca como cabeçalho o padrão dele, que é bem interessante também!
% \usepackage{fancyhdr} % Pacote para inserir cabeçalho e rodapé
% \pagestyle{fancy} % Estilo do cabeçalho e rodapé, neste caso apenas uma linha no cabeçalho
%\lhead{{\footnotesize Prandiano\textregistered}} % Inserir texto no canto esquerdo do cabeçalho
%\chead{centro} % Inserir texto no centro do cabeçalho
%\rhead{{\footnotesize Curso de Python}} % Inserir texto no canto direito do cabeçalho
%\lfoot{{\footnotesize Copyright \textcopyright 2018 Prandiano}} % Inserir texto no canto esquerdo do rodapé
%\cfoot{centro} % Inserir texto no centro do rodapé
%\rfoot{{\footnotesize Reprodução proibida: Lei 9.610 de 19/02/1998}} % Inserir texto no canto direito do rodapé
% O comando \renewcommand re-nova o parâmetro especificado
%\renewcommand{\footrulewidth}{0.4pt} % Inserir uma linha acima do rodapé
% ----------

% ----------
% Pacote para inserir múltiplas figuras:
\usepackage{subfigure} % Pacote para inserir múltiplas figuras
% \usepackage{subcaption} % Pacote incompatível com subfigure
%\usepackage{caption}
% ----------

\usepackage{multicol} % Pacote que permite usar multicolunas
\usepackage{float} % Pacote para usar o "H" do figure
\usepackage[dvipsnames]{xcolor} % Pacote de cores para textos e possivelmente outros
\usepackage{amsmath} % Pacote para usar o cases
\usepackage[makeroom]{cancel} % Pacote para usar o cancelar (em e embaixo)
\usepackage{enumitem} % Pacote para enumerar itens
\usepackage{amssymb} % Pacote para inserir os símbolos dos conjuntos numéricos N, Z, Q, R, C
\usepackage[T1]{fontenc}
\usepackage{enumitem}
\usepackage{textcomp}
\usepackage{pdfpages} % Pacote para inserir página em PDF
\usepackage{courier} % Pacote para usar a fonte courier
%\usepackage{matlab-prettifier}
\usepackage{lscape} % Pacote para mudar sentido da página
\newtheorem{theorem}{Teorema} % Pacote para inserir Teorema
\newtheorem{corollary}{Corollary} % Pacote para inserir Corolário
\newtheorem{remark}{Remark}
\newtheorem{mytheorem}{Definition}
\usepackage{hyperref} % Pacote para inserir hyperref
\usepackage{multirow} % Pacote para usar os comandos \multirow
\usepackage{pdfpages} % Pacote para incluir PDF's
\usepackage{listingsutf8} % Pacote para inserir arquivos externos tipo .txt
\usepackage{etoolbox}
\AtBeginEnvironment{verbatim}{\footnotesize}

\usepackage{standalone} % Pacote para inserir sub-arquivos LaTeX
\usepackage{bookmark} % Pacote para NÃO tratar o Índice como Apêndice e colocar letra

% Nome da listagem em português
\renewcommand{\lstlistingname}{Listagem}

% REFERÊNCIAS: ----------
%\usepackage[numbers, compress]{natbib} % Pacote para usar as referências bibliográficas
%\usepackage{usebib}
% \bibliographystyle{references} % Estilo das referências
%\setcitestyle{authoryear,aysep={},open={(},close={)}}
%\newcommand{\myand}{e}

% ----------


% PARÂMETROS DA FUNÇÃO LISTINGS: ----------
\lstset
{
    basicstyle=\small\ttfamily,
    commentstyle=\color{Green},
    keywordstyle=\color{Cerulean},
    frame=single,
    language=Python,
    morekeywords={True, False},
    numbers=left,
    numbersep=10pt,
    numberstyle=\footnotesize\color{Gray},
    showstringspaces=false,
    stringstyle=\color{Mulberry},
    tabsize=3,
    breaklines=true, % Este comando é para quebrar uma linha comprida
    postbreak=\mbox{\textcolor{red}{$\hookrightarrow$}\space}, % Este comando é para indicar com uma seta vermelha que a linha continua embaixo 
    extendedchars=true,
    literate=
        {á}{{\'a}}1 {Á}{{\'A}}1 {à}{{\`a}}1 {À}{{\`A}}1 {ã}{{\~a}}1 {Ã}{{\~A}}1 {â}{{\^a}}1 {Â}{{\^A}}1
        {ç}{{\c{c}}}1 {Ç}{{\c{C}}}1 {é}{{\'e}}1 {É}{{\'E}}1 {ê}{{\^e}}1 {Ê}{{\^E}}1 {í}{{\'\i}}1 {Í}{{\'I}}1
        {ó}{{\'o}}1 {Ó}{{\'O}}1 {ô}{{\^o}}1 {Ô}{{\^O}}1 {õ}{{\~o}}1 {Õ}{{\~O}}1 {ú}{{\'u}}1 {Ú}{{\'U}}1
        {ü}{{\"u}}1 {Ü}{{\"U}}1 {“}{{``}}1 {”}{{''}}1 {’}{{'}}1 {—}{{--}}1 {–}{{--}}1 {−}{{-}}1
        {±}{{$\pm$}}1 {σ}{{$\sigma$}}1 {→}{{$\to$}}1 {…}{{\ldots}}1
}

\lstset{inputencoding=utf8/latin1}

\lstdefinelanguage{Matlab}{
	language=Matlab,
	basicstyle=\ttfamily\small,
	keywordstyle=\color{blue}\bfseries,
	stringstyle=\color{red},
	commentstyle=\color{gray}\itshape,
	morekeywords={function, end, if, else, for, while, switch, case, otherwise},
	sensitive=true
}
% ----------

% ----------
\usepackage{nomencl} % Pacote para criar nomenclatura e lista de nomenclatura
\makenomenclature % Criar a nomenclatura
\renewcommand{\nomname}{Lista de Nomenclaturas} % Renomear o título da nomenclatura
% ----------

\usepackage{amsmath} % Pacote para inserir matrix com chaves, entre outras coisas

\usepackage{booktabs} % Pacote para usar \toprule, \midrule e \bottomrule


% EXECUTAR CÁLCULOS A PARTIR DA LINHA DE COMANDO: ---------
%\usepackage{siunitx, xfp} % Pacote para executar cálculos a partir da linha de comando
%\usepackage{l3regex}%unless you have very new expl3
%%\sisetup{locale=DE}
%%\let\oldnum\num
%%\ExplSyntaxOn
%\def\num#1{\def\tmp{#1}\regex_replace_all:nnN{,([0-9])}{.\1}{\tmp}\oldnum{\tmp}}
%\ExplSyntaxOff
% ----------


%\pagecolor{black} % Alterar a cor da página
%\color{white} % Alterar a cor do texto

% ----------
% Pacotes para inserir algoritmos e pseudo-códigos:
\setlength{\marginparwidth}{2.5cm}
\usepackage[colorinlistoftodos]{todonotes}
\usepackage{algorithm} % Pacote para inserir algoritmos
\usepackage{algpseudocode} % Pacote para inserir algoritmos e pseudo-códigos
\floatname{algorithm}{Algoritmo} % Renomear a palavra Algorithm para Procedimento
\renewcommand{\algorithmicrequire}{\textbf{entrada:}} % Renomear a palavra Input para Entrada
\renewcommand{\algorithmicensure}{\textbf{saída:}} % Renomear a palavra Output para Saída
\renewcommand{\algorithmicwhile}{\textbf{enquanto}} % Renomear a palavra While para Enquanto
\renewcommand{\algorithmicif}{\textbf{se}} % Renomear a palavra If para Se
\renewcommand{\algorithmicend}{\textbf{fim}} % Renomear a palavra End Output para Fim do
\renewcommand{\algorithmicprocedure}{\textbf{procedimento}} % Renomear a palavra Procedure para Procedimento
\renewcommand{\algorithmicfor}{\textbf{para}} % Renomear a palavra For para Para
\renewcommand{\algorithmicelse}{\textbf{senão}} % Renomear a palavra Else para Senão
\renewcommand{\algorithmicfunction}{\textbf{função}} % Renomear a palavra Function para Função 
\renewcommand{\algorithmicthen}{\textbf{então}} % Renomear a palavra Then para Então
\renewcommand{\algorithmicdo}{\textbf{fazer}} % Renomear a palavra Do para Fazer
% ----------

\newcommand{\quotes}[1]{``#1''} % Comando para inserir um texto dentro de aspas-dupla com espaço após o final

% FLUXOGRAMA: ----------
\usepackage{tikz}
\usetikzlibrary{positioning, shapes.geometric}
\tikzstyle{line}=[draw] % here

\tikzstyle{startstop} = [rectangle, rounded corners, minimum width=3cm,
minimum height=1cm,text centered, draw=black, fill=white!30]
\tikzstyle{arrow} = [thick, ->, >=stealth]
\tikzstyle{process} = [rectangle, minimum width=3cm, minimum height=1cm, text centered, draw=black, fill=white!30]
\tikzstyle{io} = [trapezium, trapezium left angle=70, trapezium right angle=110, minimum width=3cm, minimum height=1cm, text centered, draw=black, fill=white!30]
\tikzstyle{connector} = [draw, circle, fill=white!20]
\tikzstyle{decision} = [diamond, minimum width=3cm, minimum height=1cm, text centered, draw=black, fill=white!30, aspect=3]
% ----------


% CAIXA ENTORNO DAS FÓRMULAS: ----------
\usepackage{dashbox}
\newcommand\dboxed[1]{\dbox{\ensuremath{#1}}}

% ----------

%\usepackage{tabularx} % para colunas com largura ajustável

% -------- CONFIGURAÇÃO DE MARGENS --------
% \usepackage[top=2cm,bottom=2cm,left=3cm,right=2cm]{geometry}

% -------- PACOTE fancyhdr --------
\usepackage{fancyhdr}

\setlength{\marginparwidth}{2cm}

% Ajuste do tamanho do cabeçalho para acompanhar exatamente \textwidth
\setlength{\headwidth}{\textwidth}

% Define um headheight mínimo (para evitar avisos)
\setlength{\headheight}{15pt} 

% Limpa configurações padrão de cabeçalho/rodapé
\pagestyle{fancy}
\fancyhf{} 

% Marcações automáticas de seção para o cabeçalho (opcional)
\renewcommand{\sectionmark}[1]{%
	\markboth{\thesection\ #1}{}%
}

% Texto do cabeçalho à esquerda: seção e título
% Usamos \leftmark pois o \sectionmark acima grava no \leftmark
\fancyhead[L]{\small \leftmark}

% Texto do cabeçalho à direita: número de página
\fancyhead[R]{\small \thepage}

% Linha horizontal no cabeçalho (espessura 0.4pt)
\renewcommand{\headrulewidth}{0.4pt}

% Sobrescreve o título padrão “References” para “Referências”
%\renewcommand{\refname}{Referências}
% Tradução de autoref para o português
\def\figureautorefname{Figura}
\def\tableautorefname{Tabela}
\def\equationautorefname{Equa\c{c}\~ao}

% \hypersetup{backref=true}

\ifdefined\backref
        \renewcommand*{\backref}[1]{}
        \renewcommand*{\backrefentrycount}[2]{#1\ifnum#2>1 ~(#2)\fi}
\fi

\DeclarePairedDelimiter{\ceil}{\lceil}{\rceil}

\onehalfspacing

\usepackage{pgfgantt}

%++++++++++++++++++++++++++++++++++++++++++++++++++++++++++++++++++++++++++++++
% Identificacoes (se o trabalho for em inglês, insira os dados em inglês)
% Para entradas abreviadas de Professora (Profa.) em português escreva: Prof$^\textnormal{a}$.
%++++++++++++++++++++++++++++++++++++++++++++++++++++++++++++++++++++++++++++++
\course{Engenharia Aeronáutica e Mecânica} % Programa de PG ou Curso de Graduação
\area{Propulsão Aeroespacial e Energia} % Área de concentração na PG (Não utilizado no caso de TG)

% Autor do trabalho: Nome Sobrenome
\authorgender{masc} %sexo: masc ou fem
\author{\textbf{EDEN DENIS} Ferreira da Silva Lopes}{Santos}
\itaauthoraddress{Rua João Fonseca dos Santos, 20, Apto. 93}{12.230-088}{São José dos Campos--SP}

% Titulo da Tese/Dissertação
\title{Aplicação de Inteligência Artificial (AI) em Banco de Dados de Emissão de Motores}

% Orientador
\advisorgender{masc} % masc ou fem
\advisor{Prof.~Dr.}{Pedro Teixeira \textbf{LACAVA}}{ITA}

% Coorientador (Caso não haja coorientador, colocar ambas as variáveis \coadvisorgender e \coadvisor comentadas, com um % na frente)
\coadvisorgender{masc}
\coadvisor{}{}  % {Prof.~Dr.}{Carlos Alberto Rocha \textbf{PIMENTEL}}{ITA}

% Pró-reitor da Pós-graduação
\bossgender{fem}
\boss{Prof.~Dra.}{Emília \textbf{VILLANI}}

%Coordenador do curso no caso de TG
\bosscoursegender{masc}
\bosscourse{Prof.~Dr.}{ - } % masc ou fem

% Palavras-Chaves informadas pela Biblioteca -> utilizada na CIP
\kwcip{Engine taxonomy}
\kwcip{Emission classification}
\kwcip{Inteligência Artificial}
\kwcip{Artificial Intelligence}
\kwcip{Python}

% membros da banca examinadora
\examiner{Prof. Dr.}{Pedro Teixeira \textbf{LACAVA}}{President}{ITA}
\examiner{Prof. Dr.}{}{Internal Member}{ITA} % Vinícius \textbf{MALATESTA}}{Internal Member}{ITA}
\examiner{Prof. Dr.}{}{Internal Member}{ITA} % {Elisan dos Santos \textbf{MAGALHÃES}}{Internal Member}{ITA}
\examiner{Prof. Dr.}{}{Internal Member}{ITA} % {Carlos Alberto Rocha \textbf{PIMENTEL}}{External Member}{UFABC}
\examiner{Prof. Dr.}{}{Internal Member}{ITA} % {\textbf{KARL} Peter Burr}{External Member}{UFABC}

% Data da defesa (mês em maiúsculo, se trabalho em inglês, e minúsculo se trabalho em português) 
\date{14}{JULHO}{2025}

% Número CDU - (somente para TG)
\cdu{621.38}

% Folha de Registro do Documento
\newcommand{\advisorTitle}{Prof.~Dr.}
\newcommand{\advisorName}{Pedro Teixeira \textbf{LACAVA}}

\newcommand{\coadvisorTitle}{Prof.~Dr.}
\newcommand{\coadvisorName}{Carlos Alberto Rocha \textbf{PIMENTEL}}

\newcommand{\dateDay}{14}
\newcommand{\dateMonth}{JULHO}
\newcommand{\dateYear}{2025}

% Data de publicação e defesa (usada na folha de registro)
\newcommand{\frdDataPublicacao}{04 de agosto de 2025}
\newcommand{\frdNumeroRegistro}{DCTA/ITA/TD-029/2020}
\newcommand{\frdPalavrasAutor}{Big Data; Engine Emission; Emission Databank; Clustering; Artificial Intelligence}
\newcommand{\frdPalavrasResultado}{Taxonomia; Emissão de Motores; Planejamento de Dados; Inteligência Artificial; Processamento de Dados}
\newcommand{\frdApresentacao}{
	ITA, São José dos Campos. Curso de Mestrado. Programa de Pós-Graduação em Engenharia Aeronáutica e Mecânica. Área de Propulsão Aeroespacial e Energia.\\
	Orientador: \advisorTitle~\advisorName.\\
	Coorientador: \coadvisorTitle~\coadvisorName.\\
	Defesa em \dateDay/\dateMonth/\dateYear. Publicada em \frdDataPublicacao.
}

\newcommand{\PorcentagemDaLargura}{0.80}

\begin{document}

	\newpage
	\chapter{Conclusão}
		
		\section{Conclusão}
		
			\textcolor{red}{Ao desenvolver o estudo e as atividades, foi possível observar que, o projeto caminhou na direção correta, baseada em soluções históricas, a Matemática escolhida é solução para o problema de particionar as aeronaves. A aplicação do método apresentado foi desenvolvida genericamente a qual serve de ferramenta Matemática-computacional para classificar outros bancos de dados.}
			
%			Um dos grandes problemas da indústria aeroespacial, em particular a de construção de aeronaves civis e militares, é o gasto no ante-projeto desses veículos. É comum o aporte de gastos vultosos e de um grande tempo dos engenheiros projetistas para descobrir que o projeto focado é simplesmente inadequado a missão ou ao mercado. Em alguns casos, tudo isso é descoberto depois de toda a verba e tempo despendidos, o que pode ocasionar a falência da companhia. Portanto, o $ k $-means, técnica de Big Data aerotaxonométrica permite, como ficou confirmado neste trabalho, uma melhor orientação inicial para um possível projeto bem sucedido, em termos do seu início, meio e fim; evolução ajustada segundo as preferências vigentes do mercado aeronaútico. Outro sim, é quanto a evolução da Engenharia atual na direção do futuro apontado pelo Big Data que, além de prescrever a minimização do tempo de trabalho aponta também para uma redução em escala do esforço econômico-financeiro das empresas de um modo em geral.
%			
%			O $ k $-means desdobrado na dendrogramação, permite através dos seus dendros, a alocação de um futuro projeto, por semelhanças de dados ser associado pictograficamente com todos os outros projetos, já realizados, com sucesso ou fracasso seja do ponto de vista comercial, do funcional ou de ambos. Essa classificação aerotaxonômica, possibilita ao engenheiro aeroespacial, diante mão e instantaneamente, ter uma direção fundamentada para a sua futura aeronave cuja possível modificação, nessa ou naquela característica de projeto, aloca-o nesse ou naquele dendro, nesse ou no conjunto de aeronaves projetados anteriormente. Assim, durante o desenvolvimento do projeto, pode-se verificar em qualquer etapa, se o mesmo encontra-se dentro de uma classificação pré-estipulada e, caso esse tenda à um outro grupo, indesejado, pode-se fazer as modificações técnicas necessárias para que o projeto não tome outro rumo e sua verba, na maioria das vezes em \textit{déficit}, seja consumida equivocadamente. Outro ponto que toda essa Matemática destaca, são possíveis semelhanças entre aeronaves jamais comparadas devido ao grande número de projetos e de variáveis, entre outros detalhes. Não apenas isso, mas também traz a possibilidade de projetar um produto de maior custo-benefício, portanto, cada vez mais competitivo.
%			
%			Constatou-se também que o $ k $-means, numericamente funciona como uma espécie de sistema planetário onde, os medioides são as estrelas e os elementos, planetas que são gravitacionados por uma força de atração produzida por esses corpos de maior massa, medioides. Consequentemente, o modelo, ao variar a posição dos medioides (sóis) a cada iteração, consegue convergir para o ponto ótimo (de mínima ação) e, a partir daí, a posição do medioide não se altera, pois encontrou-se o ponto de equilíbrio do sistema, tudo isso sem a necessidade de se executar \quotes{infinitas} iterações já que o método converge rapidamente. Portanto, o método iterativo matemático-computacional adotado elimina os equívocos naturais das iterações anteriores, uma vez que, o critério de convergência adotado e, sobretudo, a Matemática desenvolvida, faz com que o ponto de equilíbrio seja atingido. Isto é, um elemento alocado equivocadamente em um medioide que não seja o seu ideal, será remanejado a cada passo em direção ao seu medioide ideal, até estabilizar neste.   
%			
%			No que diz respeito ao domínio comercial e, sobretudo militar, essa Matemática, atual e baseada em grandes quantidades de dados (Big Data), é ferramente fundamental na análise, comparação e desenvolvimento de projetos. Não é à toa que as empresas deste ramo da Engenharia, buscam cada vez mais profissionais com este tipo de conhecimento, que conhecem soluções não apenas específicas, mas que também permitem enxergar o mercado de maneira ampla, acima de tudo, circunstânciais. 
%			
%			Ao analisar os resultados (ver seção \ref{sec:ResultadosEDiscussoes}) chega-se a certas conclusões que somente alguém muito experiente saberia. No entanto, o $ k $-means, traz a possibilidade de uma pessoa que não seja perita neste tipo de trabalho, possa extrair conhecimento através da informação (banco de dados), mas também para quem for provecto, que este perceba certas nuances que apenas uma análise mais profunda é capaz de destacar muito além do que  já tenha feito. Vale lembrar que as discussões realizadas neste trabalho, não se limitam aqui, existe mais a ser discutido sobre tudo isso, graças aos novos métodos matemático que vem surgindo nas últimas décadas os quais trazem inúmeras possibilidades. Outro detalhe é que informações conhecidas podem ser confirmadas como é citado e comparado ao longo do texto, vale lembrar o caso do A380 e do F/A-22 Raptor discutidos na seção (\ref{sec:ResultadosEDiscussoes}).
%			
%			A Matemática é adequada e se mostrou, apesar de falibalista, com um alto grau de acerto, um excelente método para particionar bancos de dados; de fato, os resultados obtidos além de outras análises que possam ser feitas, servem de base para novos projetos e, não só isso, mas todo o estudo pode ser utilizado em outras áreas do conhecimento ou também pode ser feita diversas análises dentro do mesmo trabalho, neste caso, como se fosse uma espécie de árvores de decisão. Ao executar o particionamento do banco de dados, escolhe-se, a partir dos dendros, um grupo no qual é feita uma nova classificação que resultará em outro dendrograma e, assim, sucessivamente, pode-se extrair resultados muito mais profundos e específicos quanto às aeronaves. Da mesma forma, ao escolher um certo grupo como base de um novo projeto, esta mesma análise, pode sugerir a separação dos subsistemas  das aeronaves (superfícies aerodinâmicas, sistemas de propulsão, controle, aviônica, comunicação e componentes estruturais, dentre outros) do grupo escolhido, logo, podem ser escolhidos os subsistemas do novo projeto com o que há de melhor em termos de custo-benefício das aeronaves daquele grupo adotado inicialmente.
%			
%			O trabalho desenvolvido permite acompanhar as tendências do mercado tanto civil, quanto militar, pois é baseado, como visto, no banco de dados do que já foi produzido, ou melhor, deve-se fazer com certa frequência a análise para acompanhar tais modulações e, isto conduzirá, de fato, ao que o mercado exige em termos financeiros e ao mesmo tempo no que diz respeito à sua missão. Exemplos dessas tendências são os voos comerciais de curta distâncias, mais frequentes no mercado atual, daí a fusão entre empresas de maior porte com as de menor e as atualizações que são feitas na aviônica, em específico (mas, não apenas isso) o sistema radar dos caças que ocorrem para não se defasarem em relação à outro caça de mesma grandeza. Isso possibilita a companhia, estar sempre renovada quanto às tendências e exigências do mercado, requisito fundamental para que as empresas se mantenham na competição empresarial bem como para que uma nação mantenha sua soberania aérea. 
%			
%			Enfim, percebeu-se no $ k $-means uma possibilidade de entrever, não apenas estas, mas também outras soluções extremamente abrangentes e eminentemente funcionais para a Engenharia, porém, com um outro estilo de Matemática que contrapõe às maneiras clássicas. Fundamentado no DNA (conjunto de variáveis) das aeronaves, ao distanciar-se cada vez mais da subjetividade, a Matemática traçou as verossimilhanças entre estas, ou seja, a não absoluta, no entanto verdadeira semelhança. Este método, intercambiável, é versátil, líquido o suficiente para acompanhar as perspectivas da sociedade, logo, traz consigo, uma probabilidade real de termos em mãos a Matemática da Engenharia do futuro; no caso do presente...    

\end{document}